\section{Bridge L2 Senza Gateway}\label{bridge-l2-senza-gateway}

\subsection{shadow bridge}\label{shadow-bridge}

La topologia prevede la creazione di un bridge L2, per istanziare una
rete LAN virtuale a cui attaccare i container da verificare. I container
duplicati mantengono la topologia di rete invariata, infatti possono
mantenere gli stessi identici indirizzi IP rispetto alla rete di
produzione.

Questa topologia non prevede nessuna regola NAT e nemmeno che il bridge
abbia un'indirizzo IP, dunque a livello L3 non puo' essere visto da
nessuno.

In questo modo e' impossibile assegnare un gateway per la sottorete dei
container al bridge in quanto andrebbero in conflitto

Questa architettura e':

\begin{itemize}
\tightlist
\item
  molto semplice da implementare
\item
  efficiente: fa uso solo di strutture dati primitive del kernel linux
\end{itemize}

Tutta via rende impossibile per i container l'instradamento verso
l'esterno della rete. Dunque, dall'interno dei container di verifica,
comandi come \passthrough{\lstinline!apt!}, e
\passthrough{\lstinline!ping!} verso host esterni falliranno.

Questa architettura e' perfetta per vulnerabilità e script di verifica
che non hanno bisogno di accedere ad internet. Immaginando pero' che per
qualsiasi motivo internet potesse essere fondamentale, anche in fase di
verifica, ho scelto di scartare questa opzione.
